\documentclass[11pt,a4paper]{article}
\usepackage[utf8]{inputenc}
\usepackage[T1]{fontenc}
\usepackage{lmodern}
\usepackage{geometry}
\usepackage{hyperref}
\usepackage{xcolor}
\geometry{margin=2.4cm}
\setlength{\emergencystretch}{2em}
\newcommand{\brk}{\hspace{0pt}}

\title{\bfseries Response to the Editor and Reviewers\\[6pt]
Manuscript: \textit{Spatio-temporal wall-actuation modes and cycle-mean\\
viscous dissipation in a lid-driven cavity}}

\author{Beniaiche Ahmed, Larabi Abderrahim, Belkadi Mbarek\\
Laboratory of Fluid Mechanics, Ecole Militaire Polytechnique, Algiers, Algeria}
\date{Submission date: 21 November 2026}

\begin{document}
\maketitle

\begin{center}
\fbox{\parbox{0.92\linewidth}{\small
\textbf{Point-by-point response.} This document addresses each comment raised
by the Deputy Editor and the reviewers. For every point, we reproduce the
original comment in \textcolor{blue}{blue} and provide our response and the
corresponding revision in black. All changes are tracked in the marked-up
manuscript (\texttt{manuscript\_final\_marked.pdf}) and the marked-up
supplementary material (\texttt{supplementary\_final\_marked.pdf}).
}}
\end{center}
\medskip

\tableofcontents
\newpage

%----------------------------------------------------------------------
\section*{General Statement}
%----------------------------------------------------------------------

We thank the Deputy Editor and the reviewers for their constructive feedback.
We have completely restructured the manuscript to place the \textbf{fluid
physics} at the center of the contribution, as requested. The revised paper
now addresses the following question:

\begin{quote}
\textit{How does the spatio-temporal structure of wall actuation determine the
cycle-mean viscous dissipation of the resulting cavity flow at a prescribed
actuation-energy level?}
\end{quote}

The physics-informed neural network (PINN) is now presented as a tool for
identifying actuation branches, not as the subject of the study. The
independent multiple-relaxation-time lattice-Boltzmann (MRT-LBM) solver is
used to evaluate the hydrodynamic response and the energetic cost of each
branch. All results already published in Beniaiche (2026) have been removed
from the main text and are cited as antecedents.

The revised manuscript contains \textbf{8 sections, 8 figures, and 4 tables},
well within the 15-figure limit. The title, abstract, introduction, and
conclusions have been rewritten to lead with the actuation mode and dissipation
physics.

\vspace{1em}
\noindent\textbf{Files submitted:}
\begin{itemize}
\item \texttt{manuscript\_final.tex} and \texttt{manuscript\_final.pdf}: clean
  article file (source and author-generated PDF).
\item \texttt{manuscript\_final\_marked.tex} and
  \texttt{manuscript\_final\_marked.pdf}: marked-up version showing changes
  relative to the previous submission (old text in red, new text in blue).
\item \texttt{supplementary\_final.tex} and \texttt{supplementary\_final.pdf}:
  clean supplementary material.
\item \texttt{supplementary\_final\_marked.tex} and
  \texttt{supplementary\_final\_marked.pdf}: marked-up supplementary.
\item \texttt{RESPONSE\_TO\_EDITOR.tex} and
  \texttt{RESPONSE\_TO\_EDITOR.pdf}: this point-by-point response.
\item \texttt{figures/highlight\_image.png}: highlight image for the table of
  contents and potential cover.
\item \texttt{figures/}: source files of the figures (the 8 embedded in the
  manuscript, the 2 in the supplementary, and the earlier diagnostic figures
  kept for reproducibility).
\end{itemize}

\newpage

%----------------------------------------------------------------------
\section*{Response to the Deputy Editor}
%----------------------------------------------------------------------

\noindent\textcolor{blue}{\textbf{Editor Comment 1:} ``Physics of Fluids
publishes original theoretical, computational, and experimental work that
advances the physical understanding of fluids, meaning the mechanisms and
principles that explain how and why fluids behave as they do. As written, the
manuscript is framed as a reproducibility analysis of a physics-informed neural
network controller for the lid-driven cavity. Before it can be considered, the
paper needs to be reorganized so that the fluid physics is its central
contribution.''}

\vspace{0.5em}

\noindent\textbf{Response:} We have completely restructured the manuscript to
place the fluid physics at the center. The new title,
\textit{Spatio-temporal wall-actuation modes and cycle-mean viscous dissipation
in a lid-driven cavity}, leads with the actuation structure and the dissipation
cost, not with the neural network.

The new section organization is:
\begin{enumerate}
\item \textbf{Introduction}: presents the physical problem of wall actuation in
  a driven cavity and the central question of cycle-mean dissipation at matched
  actuation norm.
\item \textbf{Physical formulation and energetic framework}: defines the
  quadratic actuation norm $E_{\mathrm{act}}$, the viscous dissipation $D$, the
  mechanical wall power $P_{\mathrm{lid}}$, and the energy budget
  $\overline{P}_{\mathrm{lid}} = \overline{D}$.
\item \textbf{Actuation modes and hydrodynamic states}: describes the
  stationary harmonic profiles and the three identified branches B1, B2, B3 as
  distinct hydrodynamic states, not as PINN outputs.
\item \textbf{Numerical methodology}: presents the PINN as a tool for
  identifying actuation branches and the MRT-LBM as the independent solver used
  to evaluate the flow response.
\item \textbf{Energetic response of stationary wall-actuation modes}: compares
  the dissipation of uniform, $k=1$, $k=2$, and B1 mean profiles at matched
  $E_{\mathrm{act}} = 0.25$; analyzes the corner singularity and the local
  strain-rate mechanism.
\item \textbf{Cycle-mean dissipation of time-dependent actuation branches}:
  demonstrates that $\overline{D(\mathbf{u})} \neq D(\overline{\mathbf{u}})$
  and compares the cycle-mean dissipation of B1, B2, B3.
\item \textbf{Physical interpretation and implications for flow control}:
  explains why the second harmonic increases dissipation, why temporal forcing
  increases cycle-mean dissipation, and what branch accessibility tells us about
  physical performance.
\item \textbf{Conclusions}: summarizes the four main physical findings.
\end{enumerate}

The PINN branch selection, regularization-dependent state recovery, and
reproducibility audit have been moved to the supplementary material or cited as
antecedents from Beniaiche (2026). The main text now focuses on the fluid
physics: how the spatial and temporal structure of wall actuation sets the
cycle-mean viscous dissipation of the cavity flow.

\vspace{1em}

\noindent\textcolor{blue}{\textbf{Editor Comment 2:} ``The physics to bring to
the front is how the spatio-temporal structure of the wall actuation sets the
cycle-mean viscous dissipation of the cavity flow. The neural network branch
selection, the regularization dependent state recovery, and the reproducibility
audit can remain as supporting evidence, but they should not frame the work.''}

\vspace{0.5em}

\noindent\textbf{Response:} We have followed this recommendation exactly. The
new abstract begins with the physical question:

\begin{quote}
``The spatial and temporal organization of wall actuation determines how
mechanical input is distributed across the flow structures of a driven cavity,
yet the resulting viscous-dissipation cost cannot be inferred from the actuation
spectrum or from the time-mean velocity field alone.''
\end{quote}

The PINN is mentioned only as the method used to identify the actuation
branches, which are then transferred to the independent MRT-LBM solver for
evaluation. The reproducibility audit, the 185-realization campaign, the
$\lambda_{\mathrm{var}}$ sensitivity study, and the seed multiplicity analysis
have been moved to the supplementary material (Sections~S1.1--S1.4, S2.3,
S3.1--S3.3 and S5.1--S5.6).

\vspace{1em}

\noindent\textcolor{blue}{\textbf{Editor Comment 3:} ``The introduction and
reference list should be repositioned toward the fluid physics literature on
active control of the lid-driven cavity, viscous dissipation and energy budgets
in actuated cavity flows, and modal structure and state selection in driven-cavity
flow, so that the contribution is set against fluid flow.''}

\vspace{0.5em}

\noindent\textbf{Response:} The new introduction (Sec.~1) begins with the
physics of wall actuation in confined recirculating flows and cites the
canonical cavity benchmarks \cite{ghia1982high,botella1998benchmark,
shankardeshpande2000,albensoeder2005accurate,kuhlmann2019} together with the
flow-control literature carried over into the revised reference list
(Gad-el-Hak 2000; Brunton \& Noack 2015; Rowley \& Dawson 2017; Bewley
\emph{et al.}\ 2001; Bergmann \emph{et al.}\ 2005; Choi \emph{et al.}\ 1994;
Quadrio \& Ricco 2004; Corke \emph{et al.}\ 2010; Rabault \emph{et al.}\ 2019;
Garnier \emph{et al.}\ 2021), followed by the dissipation and energy-budget
literature (Pope 2000; Landau \& Lifshitz; Doering \& Gibbon; Beneddine
\emph{et al.}\ 2016). The reference list has been
reorganized to prioritize fluid physics over machine learning. The preceding
study \cite{Beniaiche2026} is discussed in one dedicated paragraph of the
introduction, which clearly distinguishes the antecedent work (branch
identification and numerical reproducibility) from the present contribution
(cycle-mean dissipation and strain-rate mechanisms); it is cited elsewhere only
where the present numbers are compared directly with that study.

\vspace{1em}

\noindent\textcolor{blue}{\textbf{Editor Comment 4:} ``Please also reduce the
number of figures to maximum of 15.''}

\vspace{0.5em}

\noindent\textbf{Response:} The revised manuscript contains \textbf{8 figures},
well below the 15-figure limit. The figures are:
\begin{enumerate}
\item Geometry and wall-actuation modes (new).
\item Spatio-temporal actuation signatures (new).
\item Grid convergence and reference sensitivity (old Fig.~10).
\item Spatial origin of dissipation differences (old Fig.~16, now central).
\item Iso-energetic comparison of the seven controls (old Fig.~11).
\item Flow topology of the actuation branches (old Fig.~13).
\item Mean velocity and fluctuation fields (old Fig.~14).
\item Phase-resolved mechanical-energy response (old Fig.~15).
\end{enumerate}

Eleven figures of the previous version are no longer in the main text. Nine of
them (old Figs.~1--9: branch-selection map, the 185-realization distribution in
modal space, spectral fingerprints, the $\lambda_{\mathrm{var}}$ probability and
entropy panels, the energy sweep, seed multiplicity, capacity/truncation and the
Fourier vs.\ Chebyshev comparison) report the identification campaigns that are
already published in Beniaiche (2026); they are now cited as antecedents, and the
selection-frequency and sensitivity numbers that the present paper still refers
to are tabulated in the supplement (Tables~S2 and S4--S9). The remaining two (the
quasi-stationarity diagnostic, old Fig.~12, and the symbolic-regression
cross-validation, old Fig.~17) are not reported at all: neither is needed for the
dissipation question posed here. The supplement contains two figures
(Figs.~S1--S2): the eight phase bins of the branch captures that underlie
Main Figs.~6--8, and the same layout with a separate stream-function scale per
panel.

\vspace{1em}

\noindent\textcolor{blue}{\textbf{Editor Comment 5:} ``Please revise the title
to lead with the actuation mode and dissipation physics.''}

\vspace{0.5em}

\noindent\textbf{Response:} The new title is:

\begin{center}
\textbf{Spatio-temporal wall-actuation modes and cycle-mean viscous dissipation
in a lid-driven cavity}
\end{center}

This title directly addresses the editor's request: it leads with the actuation
structure (``spatio-temporal wall-actuation modes'') and the dissipation physics
(``cycle-mean viscous dissipation''), and it situates the problem in the
lid-driven cavity.

\vspace{1em}

\noindent\textcolor{blue}{\textbf{Editor Comment 6:} ``Please revise your
manuscript and submit a detailed, point-by-point response to the comments. Each
response should include the original comment followed by a description of the
corresponding revision made (or a justification if no change was made). The
revised manuscript and response are due by 21-Nov-2026 via the PXP web site.''}

\vspace{0.5em}

\noindent\textbf{Response:} This document is the point-by-point response. Each
editor and reviewer comment is reproduced in blue, followed by our response and
the corresponding revision. The revised manuscript is due by 21 November 2026.

\vspace{1em}

\noindent\textcolor{blue}{\textbf{Editor Comment 7:} ``When submitting your
revision, please upload two versions of your manuscript: a clean version labeled
`Article File' and a marked-up version labeled `Marked-Up Manuscript,' which
should clearly show all changes made from the previous submission.''}

\vspace{0.5em}

\noindent\textbf{Response:} We have prepared both versions:
\begin{itemize}
\item \textbf{Clean version}: \texttt{manuscript\_final\brk.tex} and
  \texttt{manuscript\_final\brk.pdf} (26 pages).
\item \textbf{Marked-up version}: \texttt{manuscript\_final\brk\_marked\brk.tex} and
  \texttt{manuscript\_final\brk\_marked\brk.pdf} (34 pages). The marked-up version uses
  a paragraph-level diff: old text that has been removed or replaced is shown in
  \textcolor{red}{red}, new or rewritten text is shown in
  \textcolor{blue}{blue}. Obsolete figures and tables from the previous version
  are not reproduced; the corresponding changes are itemized in this response.
\end{itemize}

The same approach has been applied to the supplementary material:
\begin{itemize}
\item \textbf{Clean version}: \texttt{supplementary\_final\brk.tex} and
  \texttt{supplementary\_final\brk.pdf} (27 pages).
\item \textbf{Marked-up version}: \texttt{supplementary\_final\brk\_marked\brk.tex} and
  \texttt{supplementary\_final\brk\_marked\brk.pdf} (37 pages).
\end{itemize}

\vspace{1em}

\noindent\textcolor{blue}{\textbf{Editor Comment 8:} ``Revisions must be
submitted in the proper file formats. When submitting your revision, please
upload the source file for your revised manuscript. Figures must be embedded in
the text and visible in the PDF. An author-generated PDF must also be provided
since individual figure files are not being uploaded.''}

\vspace{0.5em}

\noindent\textbf{Response:} We have uploaded:
\begin{itemize}
\item \textbf{Source files}: clean sources \texttt{manuscript\_final\brk.tex},
  \texttt{supplementary\_final\brk.tex} (with its \texttt{data\_appendix/} table files),
  marked-up sources \texttt{manuscript\_final\brk\_marked\brk.tex} and
  \texttt{supplementary\_final\brk\_marked\brk.tex};\\
  this letter, \texttt{RESPONSE\_TO\_EDITOR\brk.tex}.
\item \textbf{Author-generated PDFs}: \texttt{manuscript\_final\brk.pdf} (26 pages),
  \texttt{supplementary\_final\brk.pdf} (27 pages), the two marked-up PDFs
  (34 and 37 pages), and
  \texttt{RESPONSE\_TO\_EDITOR\brk.pdf}.
\item \textbf{Figures}: the 8 figures of the manuscript and the 2 of the
  supplementary material are embedded in the two PDFs and visible in the text.
  Their source files (PDF and PNG) are provided in the
  \texttt{figures/} directory.
\end{itemize}

\vspace{1em}

\noindent\textcolor{blue}{\textbf{Editor Comment 9:} ``You will also be asked
to upload a highlight image that will be used in the table of contents if your
manuscript is published. The highlight image should be a visually interesting
figure representing your manuscript and may also be selected for the Cover
Image.''}

\vspace{0.5em}

\noindent\textbf{Response:} We have prepared \texttt{highlight\_image.png}
(2220\,$\times$\,900 px at 300 dpi). It shows, for the three identified
actuation branches $B_1$ (quasi-steady), $B_2$ (lid-oscillatory) and $B_3$
(mixed) at matched actuation norm $\overline{E}_{\mathrm{act}} = 0.25$ and
$Re = 500$, the spatial distribution of the cycle-mean local dissipation
density $\langle\phi\rangle$ reconstructed from the eight phase bins of the
captured periodic solution, with the streamlines of the time-mean flow
overlaid and the cycle-mean dissipation $\overline{D}$ annotated on each
panel. The three panels use a logarithmic colour scale of five decades, so
the image compares \emph{where} the dissipation is produced rather than its
absolute amplitude. This is exactly the central finding of the paper: at
matched actuation energy the three branches organise the same cavity into
radically different dissipation patterns --- space-filling two-cell
dissipation for the quasi-steady branch $B_1$, dissipation confined to the
oscillating Stokes layer under the lid for $B_2$ (whose time-mean flow
carries only $25\%$ of the cycle-mean kinetic energy), and an intermediate
organisation for $B_3$ --- and therefore different cycle-mean viscous costs. The image is generated from
the archived field captures by \texttt{scripts/make\_highlight\_image\brk.py},
so it is reproducible and uses no data other than those reported in the
manuscript. It is visually striking (three contrasting topologies on one
theme) and suitable for the table of contents and potential cover.

\newpage

%----------------------------------------------------------------------
\section*{Response to the Reviewers}
%----------------------------------------------------------------------

\noindent\textit{Note: The reviewers' comments are not reproduced here because
the Deputy Editor's comments already captured the main concerns. Points 1--14
below follow the revision plan derived from those comments; Points 15--16
additionally document the temporal mapping of the exported actuation onto the
lattice, which is the protocol item a reader of Sec.~4 needs in full; and Points
17--21 record the changes made after a further audit of the identification code
and of the discrete wall-power estimator, which is the protocol item a reader of
Secs.~2.2 and 4.1 needs in full.}

\vspace{1em}

\noindent\textbf{Point 1: New title, abstract, and introduction.}

The new title is \textit{Spatio-temporal wall-actuation modes and cycle-mean
viscous dissipation in a lid-driven cavity}. The new abstract begins with the
physical question and announces the main quantitative results (17.7\% reduction
for $k=1$, 36.8\% and 44.3\% increase for $k=2$ and B1 mean, fluctuating shares
of 92.9\%, 53.7\% and 1.8\% of the cycle-mean cost for B2, B1 and B3, and a
cycle-mean spread of the time-dependent branches by the modest factor of $1.2$).
The new introduction (Sec.~1) presents the physics of
wall actuation, the energy budget, and the central question.

\vspace{0.5em}

\noindent\textbf{Point 2: Reorganization of sections around the energy budget.}

The 8 sections follow the physical storyline: problem formulation $\to$
actuation modes $\to$ numerical methodology $\to$ stationary results $\to$
time-dependent results $\to$ physical interpretation $\to$ conclusions. The
energy budget $\overline{P}_{\mathrm{lid}} = \overline{D}$ is introduced in
Sec.~2 and used throughout.

\vspace{0.5em}

\noindent\textbf{Point 3: Elimination of repetitions with the published article.}

All results already published in Beniaiche (2026) have been removed from the
main text: second-harmonic fractions, energy sweep, 10-seed multiplicity, loss
ablation, aspect-ratio dependence, Fourier vs.\ Chebyshev, quasi-stationarity
diagnostic, same-lid-speed dissipation comparison, and symbolic regression.
The campaigns that the present argument still refers to are now reported as
supplementary tables (Tables~S2 and S4--S9) and cited as antecedents; the
quasi-stationarity diagnostic and the symbolic-regression study, which the
dissipation question posed here does not use, have been dropped from both the
manuscript and the supplement.

\vspace{0.5em}

\noindent\textbf{Point 4: Conservation of iso-energetic comparisons and
strain-rate mechanisms.}

The stationary comparison at matched $E_{\mathrm{act}} = 0.25$ (Table~2) and the
local strain-rate analysis (Fig.~4) are now central results. The corner
singularity and the reference dependence (Sec.~5.2) are preserved.

\vspace{0.5em}

\noindent\textbf{Point 5: Improvement of the power--dissipation budget.}

The energy budget (Table~4) is presented as a verification with wall-shear
estimate uncertainty, not as an exact closure. The first-order estimator gives
residuals of $-2.8\%$, $-0.6\%$ and $-0.4\%$ for B1, B2 and B3. Replacing it by a
second-order one-sided derivative moves the B1 \emph{mean-field} residual from
$-3.06\%$ to $+0.39\%$, i.e.\ across zero, which bounds the estimator
uncertainty at a few per cent, and refining the lattice from $N=256$ to $N=512$
divides the three residuals by $2.1$, $1.9$ and $2.4$ (see
supplementary Section~S4.8 and Points 15--16 below). We do not claim exact closure but document the
uncertainty.

\vspace{0.5em}

\noindent\textbf{Point 6: Reorganization of figures and tables.}

The 8 figures and 4 tables are listed in the response to Editor Comment~4. The
figure order follows the physical storyline: geometry $\to$ signatures $\to$
convergence $\to$ mechanism $\to$ iso-energetic comparison of the seven controls
$\to$ topology $\to$ mean and fluctuation fields $\to$ phase-resolved budget.

\vspace{0.5em}

\noindent\textbf{Point 7: Transfer of complementary PINN campaigns to the
supplementary.}

The 185-realization campaign, the $\lambda_{\mathrm{var}}$ sensitivity study,
the seed multiplicity analysis, and the PINN hyperparameter tables have been
moved to the supplementary material (Sections~S1.1--S1.4, S2.3 and
S3.1--S3.3; the reproducibility inventory is in Sections~S5.1--S5.6).

Additionally, the energy sweep campaign (Section~S5.1 and Table~S5 of the
version originally submitted) has
been \textbf{removed entirely} from the supplementary material. This analysis was
already documented in our previous publication (Beniaiche 2026) and is not
required for the physics-first narrative of the present manuscript. The new work
focuses on the representatives B1, B2, and B3 at $Re=500$, $E^*=0.25$, whose
characterization does not depend on the energy sweep.

\vspace{0.5em}

\noindent\textbf{Point 8: Update of the bibliography and citation of Beniaiche
(2026).}

The bibliography contains 37 references, prioritizing fluid physics over machine
learning. Beniaiche (2026) is cited in the introduction and in the
bibliography as:

\begin{quote}
A. Beniaiche, ``Physics-informed identification and independent numerical
verification of conditional modal collapse in active control of lid-driven
cavity flow,'' \textit{Phys. Fluids} \textbf{38}, 103612 (2026).
\end{quote}

\vspace{0.5em}

\noindent\textbf{Point 9: Rewriting of the conclusions and response to the
Deputy Editor.}

The new conclusions (Sec.~8) are organized around the four main physical
findings: (1) spatial structure modifies topology and strain rate; (2) ranking
depends on comparison constraint; (3) cycle-mean dissipation differs from
mean-field dissipation; (4) identified branch is not necessarily the least
dissipative. The PINN is mentioned only as a tool for identifying states.

\vspace{0.5em}

\noindent\textbf{Point 10: Explicit distinction between $E_{\mathrm{act}}$ and
$P_{\mathrm{lid}}$.}

Section~2.2 clearly states that $E_{\mathrm{act}} = \int
U_{\mathrm{lid}}^2 \, dx$ is a quadratic actuation norm, not the mechanical
work transmitted to the fluid. The mechanical wall power is $P_{\mathrm{lid}} =
\int U_{\mathrm{lid}} \tau_w \, dx$. On cycle average, $\overline{P}_{\mathrm{lid}}
= \overline{D}$.

\vspace{0.5em}

\noindent\textbf{Point 11: Addition of $U_{\max}/Re_{\max}$ column to the
stationary table.}

Table~2 now includes a column for $U_{\max}$ and $Re_{\max} = U_{\max} L / \nu$.
At $E^* = 0.25$, the uniform profile has $U_{\max} = 0.50$ ($Re_{\max} = 250$),
while the harmonics have $U_{\max} = 0.7071$ ($Re_{\max} = 354$). This
transparency avoids the objection that all flows have the same maximum Reynolds
number.

\vspace{0.5em}

\noindent\textbf{Point 12: Energy budget presented as verification with
uncertainty.}

Section~6.5 and Table~4 present the energy budget as a verification with
wall-shear estimate uncertainty. The recovered-velocity, first-order estimator
gives cycle-mean residuals $\overline{P}_{\mathrm{lid}}^{\mathrm{rec}}-\overline{D}$ of
$-2.8\%$, $-0.6\%$ and $-0.4\%$ for $B_1$, $B_2$ and $B_3$, i.e.\ the budget
closes to a few per cent. The estimator built on the \emph{prescribed} lid
velocity gives $+3.4$, $+4.6$ and $+4.7\%$, and the two brackets bound the
estimator uncertainty (Section~S4.8). We do not claim exact closure. All of these
percentages, together with the alternative estimators, are recomputed from the
archived field captures by
\texttt{scripts/d2\_derived\_quantities\brk.py}, whose output is
\texttt{results/d2\_derived/d2\_summary.json}; the residual annotated in Fig.~8 is
the one of the two recorded signals ($-2.8$, $-0.6$ and $-0.4\%$), and the caption
states which estimator each number comes from.

\vspace{0.5em}

\noindent\textbf{Point 13: Marked-up versions built by sharing the new preamble.}

The marked-up manuscript and supplementary were built using a paragraph-level
diff script (\texttt{mark\_manuscript.py}) that shares the new preamble for both
the old and new bodies. This ensures that the marked-up files compile without
errors related to obsolete macros or missing packages.

\vspace{0.5em}

\noindent\textbf{Point 14: Deep point-by-point verification.}

We have performed a deep verification of the revision plan and the editor's
comments. The reproducibility inventory and the data-availability statement are
provided in the supplementary material
(Sections~S5.5--S5.6). All points have been addressed.

\vspace{0.5em}

\noindent\textbf{Point 15: Temporal mapping of the exported actuation onto the
lattice.}

The protocol is stated in Sec.~4.2 of the manuscript and verified in
supplementary Sections~S4.3 and S4.5. The actuation is defined over the physical
nondimensional period $T=2$, with Fourier components
$\cos(2n\pi t/T)=\cos(n\pi t)$, so the fundamental actuation frequency is
$\omega_1=2\pi/T=\pi$ in cavity units on every grid. The lattice advances by
$\Delta t=\Delta x/U_{\mathrm{ref}}$ with $\Delta x=1/N$ and
$U_{\mathrm{ref}}=0.05$; one physical period therefore spans
$N_T=TN/U_{\mathrm{ref}}=40N$ lattice steps, i.e.\ $N_T=10\,240$ at $N=256$ and
$20\,480$ at $N=512$ (Eq.~14 of the manuscript). Each of the $256$ exported
samples of the waveform is held for $N_T/256$ consecutive steps, i.e.\ $40$ at
$N=256$ and $80$ at $N=512$, so its physical duration $T/256=7.8125\times10^{-3}$
is resolution-independent and the staircase it imposes is common to both grids.
The largest instantaneous lid speed of the campaign is $1.23\,U_{\mathrm{ref}}$,
hence $\mathrm{Ma}\approx0.11$, and the oscillatory wall-layer thickness
$\delta_1=\sqrt{2\nu/\omega_1}=0.036$ is about nine cells at $N=256$ and eighteen
at $N=512$, so the same margin resolves that layer at both resolutions.

Three checks close the mapping. \emph{(i) Frequency response.} A dedicated run
whose lid is the pure harmonic $\sin(2\pi x)\cos(\omega_1 t)$, at the same
$E^\ast$, $Re$, $N$ and $\omega_1$
(\texttt{scripts/validate\_stokes\_analytic\brk.py}, Section~S4.3) records at the
antinode a single-sided amplitude of $0.927$ at $\omega_1$ and no other
significant line --- the imposed unit amplitude, recovered to within the
$7.3\%$ slip of the velocity boundary condition --- and a wall-normal exponential
decay of fitted length $5.4$ cells ($R^2=0.999$ over the eleven cells nearest the
wall) against the half-space value $\delta=\sqrt{2\nu/\omega_1}=9.1$ cells. The
layer is of the expected order and $41\%$ thinner, which is attributed to the
non-uniform lid, the finite depth (about $28\delta$) and the steady shear
superimposed by the return flow; the run is used as a frequency-response check of
the solver and of the mapping, not as a quantitative benchmark of the layer
thickness. \emph{(ii) Spectral purity of the branch runs.} A streamwise probe at
the cavity centre, recorded over the whole sampling window, has its spectral peak
at $3.141593$ for all three branches, with $97.5$--$99.5\%$ of its fluctuation
energy in the five spectral bins around that peak, and the dissipation signal
peaks at $\omega_1$ for the quasi-stationary $B_1$ and at $2\omega_1$ for $B_2$
and $B_3$ (Section~S4.5). \emph{(iii) Effect of the hold.} Replaying $B_2$ at
$N=256$ with the exact sample-time command instead of the $40$-step hold returns
$\overline{D}=1.4425\times10^{-2}$ against $1.4408\times10^{-2}$, i.e.\
$+0.114\%$, with the fluctuating part moved by $+0.12\%$ and the kinetic energy
by $+0.005\%$, while the realized cycle-mean actuation energy stays within
$1\times10^{-16}$ of the target $E^\ast=0.25$ on both grids; the
staircase is therefore not the source of the reported differences.

Every cycle mean quoted in the paper is a direct time average over the $32$
retained forced periods, after a transient terminated by a drift test
($24$ forced periods, successive three-period windows changing by less than
$2\times10^{-3}$). The eight-bin phase quadrature of the same data reproduces
$\overline{D}$ to within $0.004\%$ but under-estimates the peakier lid power by
$0.14\%$, $5.5\%$ and $3.9\%$ for $B_1$, $B_2$ and $B_3$, so the quadrature is
used only as a cross-check.

\vspace{0.5em}

\noindent\textbf{Point 16: Grid-refinement sensitivity of the time-dependent
comparison.}

The three branch runs were replicated at $N=512$ under the same temporal
protocol, and the comparison is reported as a \emph{grid-refinement
sensitivity}: two grids bound the variation of the cycle means, they do not
determine a convergence order, so no extrapolation and no convergence index is
computed for them (Sec.~4.3 and Table~S15). The refined runs raise
$\overline{D}$ by $5.4\%$, $7.3\%$ and $7.1\%$ for $B_1$, $B_2$ and $B_3$,
preserve the ordering $B_2>B_1\simeq B_3$, and narrow the $B_1$--$B_3$ gap from
$2.1\%$ to $0.5\%$, so the central conclusion of Sec.~6 --- that the three
branches lie within a factor $1.23$ of one another on the cycle average while
their fluctuating parts span two orders of magnitude --- is unchanged by the
refinement. The power-balance residual of Sec.~6.5 is divided by $2.1$, $1.9$
and $2.4$, i.e.\ the factor of two a first-order one-sided stencil predicts under
$\Delta x\to\Delta x/2$, and moves to $-1.4\%$, $-0.3\%$ and $-0.2\%$, while the
largest slip of the velocity boundary condition drops by $38$--$47\%$. The
periodicity of the averaged state is not degraded: the half-window phase
reproduction of the finer captures is $1.79\times10^{-3}$,
$1.48\times10^{-5}$ and $2.69\times10^{-4}$ against $1.87\times10^{-3}$,
$1.46\times10^{-5}$ and $2.61\times10^{-4}$ at $N=256$. Each finer run cost
$\approx13.1$~h of single-threaded wall time against $\approx0.7$~h for the
corresponding $N=256$ capture, which is why the phase-resolved fields are
reported on the $N=256$ realization and the cycle-mean scalars on both grids.
Sec.~4.3 and Supplementary Sec.~S4.5 now state that division of labour
explicitly --- the replication supplies scalar diagnostics only, and every field
and phase-resolved quantity in the paper is an $N=256$ capture --- so that no
figure is read as finer than it is.

\vspace{1em}

\noindent\textbf{Point 17: Two time variables in the identification, now
distinguished.}

A fresh audit of the training script confirms that the identification runs on a
raw nondimensional clock $\tau$ sampled uniformly over $[0,T_{\max}]$ with
$T_{\max}=1$, that the control is expanded there as
$\sin(m\pi x)\cos(n\pi\tau)$, and that the momentum residual differentiates with
respect to $\tau$ directly, with no $1/T$ chain factor. The manuscript and the
supplement previously described that same variable as the normalized phase
$\theta=t_{\mathrm{phys}}/T$ of the exported period, which conflated the two
clocks. Sec.~4.1 and Supplementary Secs.~S1.1, S1.2 and S1.4 now state the
distinction explicitly, in the order the computation performs it: the coefficient
network returns the amplitudes $a_{mn}(Re)$, which carry no time variable; the
flow network is trained on $\tau\in[0,1]$, which covers one half-period of the
fundamental; and the exported wall actuation is the \emph{same} Fourier series
re-evaluated on the physical clock over its complete period,
$\cos(2n\pi t_{\mathrm{phys}}/T)=\cos(n\pi t_{\mathrm{phys}})$ with
$0\le t_{\mathrm{phys}}\le T=2$, so the fundamental of what the lattice actually
sees is $\omega_1=2\pi/T=\pi$ (Sec.~4.2, Eq.~14, Supplementary Sec.~S2.1). The PINN
flow field is not periodicised and not transferred to the solver --- no endpoint
condition is imposed on it --- and the periodic hydrodynamic response is
established independently by the MRT-LBM calculation and verified a posteriori to
$10^{-3}$ in relative field norm (Supplementary Sec.~S1.4). We also state why the
constraint is not affected by the shorter training window:
$\int_0^1\cos^2(n\pi\tau)\,d\tau=\int_0^1\cos^2(2n\pi\theta)\,d\theta=1/2$ with
vanishing cross-terms on both windows, so the penalised quantity is the cycle
mean of the exported signal even though the flow field was trained on half a
period (Supplementary Sec.~S2.2).

\vspace{1em}

\noindent\textbf{Point 18: Direction of the Stokes comparison.}

The abstract, the contribution list of Sec.~1, Sec.~6.3 and Sec.~8 previously said
that the second-Stokes estimate ``accounts for $95\%$ and $98\%$'' of the
fluctuating dissipation. That is the reciprocal of what the numbers say: the
measured fluctuating parts are $1.339\times10^{-2}$ for $B_2$ and
$6.278\times10^{-3}$ for $B_3$ against estimates of $1.411\times10^{-2}$ and
$6.383\times10^{-3}$, so the measurements are $95\%$ and $98\%$ \emph{of the
estimate} and the estimate is slightly the larger. The four places now read that
``the estimate agrees with the measured fluctuating dissipation to within
approximately $5\%$ for $B_2$ and $2\%$ for $B_3$, without an adjustable
constant'', which states the accuracy of the approximation rather than a share of
the budget. For the same reason the abstract no longer claims the estimate
``locates'' the cost in the wall layer: the model is a half-space one, so the
agreement supports a wall-layer interpretation of the dominant fluctuating cost
rather than a spatial decomposition of the total, and Sec.~7.1 adds explicitly
that the $56\%$ the estimate leaves unexplained for the $k=1$ profile is not
thereby assigned to the cavity core.

\vspace{1em}

\noindent\textbf{Point 19: Definition of \(R_K\).}

$R_K$ was described in the nomenclature as the ``ratio of fluctuating to mean
kinetic energy'', but the reported values are computed against the total. With
$\overline K=\overline K_{\mathrm{mean}}+\overline K_{\mathrm{fluct}}$ the
definition is now
$R_K=\overline K_{\mathrm{fluct}}/\overline K$, ``the fraction of the cycle-mean
kinetic energy carried by temporal fluctuations'', which is what reproduces the
table: $B_1$ $0.16\%$ against a mean field holding $99.8\%$, $B_2$ $74.9\%$
against $25\%$, $B_3$ $12.5\%$ against $87.5\%$. The two supplement statements
that read ``fluctuates by $R_K$ of its mean'' in Secs.~S4.3 and S4.5 now read
``of the cycle-mean kinetic energy'', which matters most for $B_2$, where the old
wording would have implied that the mean field retains three quarters of the
energy instead of a quarter.

\vspace{1em}

\noindent\textbf{Point 20: Geometry of the discrete wall-power estimator.}

The main text and the supplement disagreed about where
$\overline{P}_{\mathrm{lid}}^{\mathrm{rec}}$ is evaluated: Sec.~2.2 placed both
factors of the product on the plane $y=1-\Delta y$ (index $N-2$) while
Supplementary Sec.~S4.8 wrote them on row $N-1$. We re-read the estimator in the
solver, and it uses the lid node row for both factors,
$U_{\mathrm{wall}}^{\mathrm{rec}}(x_j)=u_x(N-1,j)$ and
$\tau_w^{\mathrm{FD}}(x_j)=Re^{-1}[u_x(N-1,j)-u_x(N-2,j)]/\Delta y$, with
trapezoidal end weights. Eq.~(5) now carries those discrete definitions and
Sec.~2.2, the caption of Table~4, Sec.~6.5 and Supplementary Sec.~S4.8 use the
same single convention. The accompanying statement about the boundary condition
was corrected in the same pass: the lid is imposed by half-way bounce-back with
momentum injection into the two wall-facing populations, which enforces the moving
wall at the hydrodynamic boundary half a lattice spacing outside row $N-1$; it
therefore does \emph{not} make the recovered node velocity equal the commanded
profile, and Supplementary Sec.~S4.1 no longer claims it does. That discrepancy is
the slip already quantified in Sec.~6.5 --- $6.0\%$, $11.5\%$ and $11.8\%$ of the
cavity speed at $N=256$, dropping to $3.2\%$, $7.2\%$ and $6.9\%$ at $N=512$ ---
and Sec.~2.2 now explains the residuals of Table~4 through it, instead of leaving
the reader to infer why the recovered and prescribed estimators bracket the
closure.

\vspace{1em}

\noindent\textbf{Point 21: Quantitative statements made exact.}

Four imprecise phrases were replaced by the numbers they summarise. Sec.~6.2, the
caption of Fig.~5 and Sec.~8 said the branch totals lie ``within twenty per cent''
of one another; $B_2$ is in fact $20.7\%$ above $B_1$ and $23.2\%$ above $B_3$ at
$N=256$, and Sec.~6.2 and the Conclusions now give both ratios and bound the
spread by a factor $1.23$. ``Indistinguishable within $2\%$'' for $B_1$ and $B_3$
became ``nearly equal, $2.1\%$ apart at $N=256$ and $0.5\%$ apart at $N=512$'' in
the abstract, Sec.~6.2, Sec.~8 and Supplementary Sec.~S4.5, which keeps the claim
true at both resolutions instead of rounding one of them. The nomenclature now
distinguishes $U_0$, the scaling velocity of the nondimensional problem, from
$U_{\mathrm{ref}}=0.05$, the lattice reference velocity that sets
$\Delta t=U_{\mathrm{ref}}/N$, and lists $E_{\mathrm{act}}$ as dimensionless with
dimensional equivalent $U_0^2L$, since the constraint is written in the scaled
variables of Eq.~(1). Finally, the single-harmonic test at $n=1,2,4$ is stated as
work not yet done in Sec.~6.3 and Sec.~8: Supplementary Sec.~S4.5 tabulates the
harmonic budgets of the three controls already used, it does not contain a
mono-harmonic campaign, and the sentence pointing to it no longer suggests
otherwise.

\newpage

%----------------------------------------------------------------------
\section*{Summary of Changes}
%----------------------------------------------------------------------

\noindent\textbf{Title:} Changed from ``Physics-informed identification and
energetic characterization of conditional modal collapse in active control of
lid-driven cavity flow'' to ``Spatio-temporal wall-actuation modes and
cycle-mean viscous dissipation in a lid-driven cavity.''

\noindent\textbf{Abstract:} Completely rewritten to lead with the physical
question and announce the main quantitative results.

\noindent\textbf{Introduction:} Completely rewritten to present the physics of
wall actuation, the energy budget, and the central question. The PINN is
mentioned only as a tool for identifying branches.

\noindent\textbf{Sections:} Reorganized from 9 sections (focused on PINN
identification) to 8 sections (focused on fluid physics); the symbolic-regression
section of the previous version has been dropped.

\noindent\textbf{Figures:} Reduced from 17 figures to 8 figures. Six figures are
carried over (old Figs.~10--11, 13--16), two are new, and the eleven dropped
figures are the nine identification-campaign figures already published in
Beniaiche (2026) plus the quasi-stationarity and symbolic-regression diagnostics.

\noindent\textbf{Tables:} Reduced from 11 tables to 4 tables in the main text. The
detail tables --- those already published in Beniaiche (2026) and those reporting
the identification hyperparameters, the seed-level metrics and the
grid-convergence data --- are collected in the supplement, which now holds
19 tables (Tables~S1--S19).

\noindent\textbf{References:} Reorganized to prioritize fluid physics over
machine learning. Added Beniaiche (2026) as an antecedent. The bibliography grew
from 25 to 37 entries, each checked against the publisher record before citing.

\noindent\textbf{Supplementary material:} Restructured into five sections
(S1--S5): PINN implementation, control-space and modal-analysis details,
identification-campaign robustness, the independent MRT-LBM verification with
grid uncertainty, and the reproducibility inventory. It contains 19 tables,
2 figures and its own short bibliography.

\noindent\textbf{Marked-up versions:} Prepared for both the manuscript and the
supplementary material using a paragraph-level diff script.

\noindent\textbf{Highlight image:} Prepared showing the strain-rate distribution
for the three actuation branches.

\newpage

%----------------------------------------------------------------------
\bibliographystyle{plain}
\begin{thebibliography}{10}

\bibitem{ghia1982high}
U.~Ghia, K.~N.~Ghia, and C.~T.~Shin,
``High-{Re} solutions for incompressible flow using the {Navier--Stokes}
equations and a multigrid method,''
\textit{J. Comput. Phys.} \textbf{48}, 382--411 (1982).

\bibitem{botella1998benchmark}
O.~Botella and R.~Peyret,
``Benchmark spectral results on the lid-driven cavity flow,''
\textit{Comput. Fluids} \textbf{27}, 421--433 (1998).

\bibitem{shankardeshpande2000}
S.~Shankar and V.~Deshpande,
``Numerical simulation of two-dimensional lid-driven cavity flow at high
{Reynolds} numbers,''
\textit{Int. J. Numer. Methods Fluids} \textbf{32}, 11--29 (2000).

\bibitem{albensoeder2005accurate}
S.~Albensoeder and H.~C.~Kuhlmann,
``Accurate three-dimensional lid-driven cavity flow,''
\textit{J. Comput. Phys.} \textbf{206}, 536--558 (2005).

\bibitem{kuhlmann2019}
H.~C.~Kuhlmann and F.~Romanò,
``The lid-driven cavity,''
in \textit{Open Problems in Fluid Mechanics}, edited by H.~C.~Kuhlmann and
F.~Romanò (Springer, 2019), pp.~233--309.

\bibitem{Beniaiche2026}
A.~Beniaiche,
``Physics-informed identification and independent numerical verification of
conditional modal collapse in active control of lid-driven cavity flow,''
\textit{Phys. Fluids} \textbf{38}, 103612 (2026).

\end{thebibliography}

\end{document}
